\section*{Mission Control narration: how to use these scripts}
\label{narration-start}
\begin{panel}{Original game dialogue --- not historical quotations}
Everything in quotation marks in this appendix is invented read-aloud
dialogue or scene description. It is not a NASA transcript.
Use it verbatim or adapt it to your crew, without changing the inputs,
grading bands, resource rules, or historical event anchors.
\end{panel}

\textbf{Before a card.} Read its briefing below, announce the GET
printed on the calculation card, and read that card's telemetry and orders.
Where the card supplies no exact event time, name the phase instead;
do not invent a new deadline. Briefings intentionally omit answers.

\textbf{After commitment.} Grade first, subtract the appropriate
command-reserve charge once, and check the new reserve.
Read only the matching GREEN, AMBER, or RED example.
For an error, the examples below assume reserve is still nonnegative.
If reserve is negative, use the appropriate ending instead of inviting
another attempt. A correct sheet never receives a scripted error.

\textbf{Correction and continuation.} GREEN narration authorizes the
nominal game action. AMBER and RED narration hold that action pending a
corrected pad. After correction, use the common release line below,
then the card's GREEN scene where appropriate. For 11-06 AMBER, retain
the less desirable landing site; for 11-07 retain the recorded science
flag rather than announcing full science.

\begin{panel}{Reusable referee lines}
\textbf{Commitment:} ``Flight crew, give me your signed command,
units, and GO or NO-GO decision. Is that your committed pad?''

\textbf{AMBER bookkeeping:} ``That is an AMBER checkpoint.
One command-reserve point is charged. Your reserve is now
[read the ledger]. Revise the pad before we proceed.''

\textbf{RED bookkeeping:} ``That is a RED checkpoint.
Two command-reserve points are charged, once for this sheet.
Your reserve is now [read the ledger].''

\textbf{Release after correction:} ``The corrected pad is accepted.
There is no additional charge for this revision. Resume from this
card's nominal calculation baseline; keep the warning on the ledger.''

\textbf{Unmodeled passage of time:} ``We now advance to the next
card's supplied epoch. The intervening routine work is narrated;
no extra consumables are deducted.''
\end{panel}

\textbf{Predictions are not actual disasters.} Error scenes describe
the simulator's warning or contingency rehearsal. Do not narrate a
collision, exhausted tank, or atmospheric entry as having occurred
while players are still allowed to correct their commands.
No scene adds dice rolls, real-time pressure, hidden fuel charges, or
extra calculations. The cartridge scene is dramatic description,
not a breathing-apparatus construction procedure.

\clearpage
\section*{Apollo 11: launch and narrated transitions}
\textbf{Opening / before 11-01.}
``Armstrong, Aldrin, Collins: your mission is to land Eagle, complete
the planned surface work, and bring the crew home. The Saturn rises
through the cloud layer. Staging is nominal; Earth curves beneath you.
At planned GET 000:11:43, the spacecraft reaches its parking orbit.
For this game the launch is successful. Houston is preparing your
lunar-departure pad. Take out your pencils and slide rule.''

\textbf{After corrected 11-01 / before 11-02: docking and extraction.}
``The third-stage engine shuts down and the cabin settles into
weightlessness. Collins turns Columbia to face Eagle, eases toward
the docking target, and secures the connection. The lunar module is
extracted from the stage. That routine docking is narrated here,
not a new calculation checkpoint. Now the spent stage must be
left safely behind: prepare the avoidance pad.''

\textbf{After corrected 11-02 / before 11-03: translunar coast.}
``The spent stage recedes. Earth becomes a smaller blue disk while
the Moon grows ahead. The crew alternates navigation observations,
housekeeping, meals, and rest. Routine midcourse updates and thermal
control are handled in the narrated coast. Houston now sends the
lunar-capture pad. Check it before the Moon cuts off our radio link.''

\textbf{After corrected 11-03 / before 11-04: orbit operations.}
``Radio contact returns as Columbia comes around the Moon.
For the game's accepted capture branch, you are in lunar orbit.
The second orbit-insertion maneuver and routine LM inspection are
completed as narrated flight-plan operations. Collins has daylight
observations and communications to schedule; these are not the
same kind of window.''

\textbf{After corrected 11-04 / before 11-05: undock.}
``Armstrong and Aldrin complete Eagle's checkout.
Collins remains in Columbia, monitoring the separation.
The docking connection releases and the two vehicles drift apart.
Separation and activation are narrated operations here.
Eagle's next checkpoint is the burn that lowers its perilune.''

\textbf{After corrected 11-05 / before 11-06: descent readiness.}
``Eagle coasts toward the powered-descent point. Aldrin reviews the
corrected timing while Armstrong watches the surface move below.
Houston calls for the final descent-readiness check. No new coast
budget is charged. At the next card's terminal exercise, decide how
much diversion the supplied fuel model permits.''

\textbf{After corrected 11-06 / before 11-07: surface operations.}
``Eagle is settled on the surface in the accepted landing branch.
The crew completes the contact checks, secures the spacecraft,
and prepares for the planned walk. Collins continues his orbiting
watch. Before either walker leaves, Houston wants the oxygen
and science schedule checked.''

\textbf{After corrected 11-07 / before 11-08: prepare ascent.}
``The walkers return and secure the cabin. Record the science result
already earned; do not restore curtailed work.
Equipment stowage, rest, and ascent preparations are narrated.
Collins is ready in Columbia. Eagle's crew must now check ascent
timing and the separate close-range braking exercise.''

\textbf{After corrected 11-08 / before 11-09: jettison.}
``The corrected approach brings the docking targets together.
Contact, capture, and a secure connection are confirmed.
Armstrong and Aldrin transfer back to Columbia.
The LM-jettison sequence and separation maneuver are narrated
without another calculation. All three crew now prepare
the trans-Earth injection pad.''

\clearpage
\section*{Apollo 13: launch, crisis, and narrated transitions}
\textbf{Opening / before 13-01.}
``Lovell, Haise, Swigert: the original mission is a lunar landing.
Launch, Earth-orbit checkout, lunar departure, and docking are
successful for this game. Odyssey and Aquarius head outward together.
Routine coast operations pass without a calculation checkpoint.
Now announce GET 055:54:53: a sharp disturbance runs through the
spacecraft, and warning indications appear on the panels.
Something has failed in the service module. The landing mission is
over; your new mission is to bring three people home.''

\textbf{After corrected 13-01 / before 13-02: transfer.}
``Odyssey's nonessential loads go dark.
Haise brings Aquarius into its lifeboat role while Swigert preserves
the command module for later. The crew moves through the connecting
tunnel with checklists and essential equipment.
Houston now needs a return maneuver that the LM engine can perform
with the two spacecraft still docked.''

\textbf{After corrected 13-02 / before 13-03--04: planning audit.}
``The accepted correction puts the simulator on its return branch.
Houston and the crew spread out the consumable sheets.
The oxygen and power audits both start their accounting at
058:00:00, even though this discussion follows the maneuver scene.
They cover the whole remaining trip; they are not fresh snapshots
or deductions from an already-used inventory.''

\textbf{After corrected 13-04 / before 13-05: lunar swingby.}
``Aquarius is quiet and dim under the accepted conservation plan.
The Moon fills the windows and then passes behind the docked stack.
For a time there is no radio contact. When Houston's voice returns,
the crew prepares PC+2: the major burn intended to shorten the trip
and improve the return target. Read the mass model carefully.''

\textbf{After corrected 13-05 / before 13-06: return coast.}
``The burn ends and Earth becomes the destination ahead.
The crew settles back into the low-power return coast.
Houston now reviews cooling and drinking water.
This is another audit of the same full horizon starting at 058:00:00,
not an extra charge for the time just narrated.''

\textbf{At 13-07: cartridge-adapter assembly.}
Read the briefing and telemetry, then grade the committed plan.
Only after the plan is accepted or corrected, read:
``Haise lays out the spare cartridges and the materials available
in the cabin. Swigert reads the simulated ground instructions while
Lovell holds the parts steady in weightlessness.
The crew assembles the adapter during the card's stated assembly
interval. For this game, the adapter connects the allocated
cartridges to the airflow path and the capacity plan is accepted.
No extra crafting roll or chemical calculation is required.''

\textbf{After corrected 13-07 / before 13-08--09: manual navigation.}
``The adapted scrubber arrangement is working in the accepted game
branch. With the ordinary guidance setup unavailable for this
exercise, Lovell uses the visible Earth as an attitude reference.
Haise reads the signed correction pad, and Swigert checks the units.
The later entry predictor will ask a different question:
how long must the final correction last?''

\textbf{After corrected 13-09 / before 13-10: ready Odyssey.}
``The corrected entry pad is accepted.
The damaged service module is separated in the narrated sequence.
Swigert prepares to wake Odyssey while Aquarius is still available.
Do not jettison the lifeboat until the final card's energy and
reactivation timing gates have been accepted.''

\newcommand{\NarrationCard}[6]{%
  \subsection*{#1 --- #2}
  \textbf{Briefing (before calculation):} ``#3''\par
  \textbf{GREEN (accepted command):} ``#4''\par
  \textbf{AMBER (charge 1, then correct):} ``#5''\par
  \textbf{RED (charge 2; if reserve permits, then correct):} ``#6''\par}

\clearpage
\section*{Apollo 11 checkpoint scripts / 11-01--02}
\NarrationCard{11-01}{Commit to the Moon}
{Houston to Columbia: we are at the lunar-departure gate.
Aldrin, read the velocity units. Collins, check the duration and
decimal place. Armstrong, hold the command until the crew commits.}
{The departure pad is accepted. The third stage lights, acceleration
builds, and the burn ends on the checked clock.
Columbia and Eagle begin the game's outbound lunar coast.}
{Guidance review has flagged a targeting discrepancy.
We need a corrective update before accepting the departure pad.
Check your unit conversion, average acceleration, and cutoff time.}
{The simulator rejects this departure solution.
Run the abort-and-replan rehearsal; do not execute this pad.
If command reserve permits, rebuild the calculation and read it back.}

\NarrationCard{11-02}{Clear the spent stage}
{Columbia, the docked spacecraft is clear of the extraction sequence,
but the stage is still nearby. Read the vector as one maneuver,
then tell Houston its magnitude and commanded duration.}
{The safe-separation predictor accepts your pad.
The short maneuver sends the spent stage steadily away.
The crew can turn its attention to the coast ahead.}
{The predicted separation rate differs from the accepted model.
Houston requests a tracking update and corrected magnitude.
Keep the component directions on the revised pad.}
{The safe-separation check fails. Rehearse emergency avoidance in
the simulator; no collision has occurred in the game.
Rebuild the vector and duration before proceeding.}

\clearpage
\section*{Apollo 11 checkpoint scripts / 11-03--04}
\NarrationCard{11-03}{Capture lunar orbit}
{The Moon is about to block Houston's signal.
Collins, verify the retrograde signs; Aldrin, check the magnitude
and cutoff clock. Your committed pad must stand on its own.}
{The engine shuts down at the accepted cutoff.
As the spacecraft emerges from behind the Moon, radio contact
returns. The simulator accepts lunar capture.}
{Tracking predicts an off-target lunar orbit.
Houston requests a trim update. Correct the signed pad and timing
before we accept the nominal capture branch.}
{The simulator rejects lunar capture with this pad.
Prepare the contingency-return rehearsal.
If reserve remains, correct the burn solution before resuming.}

\NarrationCard{11-04}{Work the lunar clock}
{Collins, sunlight and radio coverage are separate constraints.
Houston needs the observation schedule checked against daylight,
not against a coverage duration that happens to look similar.}
{The daylight task fits, and the schedule is accepted.
Record the coverage fraction separately; we have not assumed
that the entire observation has continuous radio contact.}
{The observation timing would lose useful coverage.
Houston requests a revised schedule. Recheck the orbit period,
daylight interval, and coverage percentage independently.}
{The handover schedule is rejected.
Rebuild the scheduling sheet; do not substitute a radio window
for a daylight window. The missed handover is a simulator contingency.}

\clearpage
\section*{Apollo 11 checkpoint scripts / 11-05--06}
\NarrationCard{11-05}{Lower Eagle's perilune}
{Eagle, Houston wants descent-orbit insertion checked before ignition.
Use the required velocity change specified on this card, then tell
us when cutoff occurs and how much coast remains before PDI.}
{The DOI pad is accepted.
The burn ends on the corrected clock, and Eagle enters the
game's descent coast with its readiness interval accounted for.}
{The coast briefing would be late on this sheet.
Expedite the pad review in the narrative, then correct the cutoff
and coast interval. No extra fuel or elapsed-time charge is added.}
{The descent-readiness check fails.
Rehearse the full readiness sequence before attempting PDI.
If reserve remains, correct the sheet and resume the nominal coast.}

\NarrationCard{11-06}{Land or wave off}
{Armstrong sees a rough patch ahead; Aldrin watches altitude and fuel.
Check the supplied terminal model and choose the longest listed
diversion that still permits descent without exhausting the fuel.}
{The chosen diversion passes the fuel gate.
Eagle moves toward the selected patch, descends, and settles.
The crew reports contact and secures the spacecraft.}
{The landing pad needs a last-second correction.
Correct the calculation before touchdown is accepted.
In the accepted continuation, Eagle lands at a less desirable site;
the warning remains on the ledger.}
{The fuel predictor rejects the landing command.
Rehearse an immediate wave-off; the game has not committed touchdown.
If reserve remains, correct the diversion choice and landing pad.}

\clearpage
\section*{Apollo 11 checkpoint scripts / 11-07--08}
\NarrationCard{11-07}{A walk with a turnaround time}
{The hatch is ready, but Houston wants the return margin protected.
Budget each walker's oxygen separately, then check how many
complete science stations fit the planned EVA.}
{The EVA plan is accepted. The walkers work through the planned
science schedule and return with the protected margin intact.
Record full planned science if no warning occurred on this card.}
{Houston calls for an early turnaround.
Record reduced science, then correct the oxygen and station sheet.
The crew returns safely in the accepted continuation, but the
lost science is not restored by the correction.}
{The plan violates the modeled return margin.
Record EVA curtailed and rehearse the immediate return.
Correct the sheet before ascent preparations; do not claim that
the original full science schedule was completed.}

\NarrationCard{11-08}{Rise, meet, and dock}
{Eagle's ascent pad is on the desk.
Check insertion timing first. Then work the separate close-range
exercise: how long before contact, and which way must you brake?}
{Ascent timing and braking are accepted.
Eagle reaches the game's rendezvous branch; the closing motion
is canceled. The docking targets align, contact is made,
and the connection is secured.}
{Collins holds the docking approach while the crew corrects its pad.
Check the insertion clock and braking duration.
The narrative delay adds no hidden propellant charge.}
{The docking predictor flags a possible collision.
Rehearse emergency separation in the simulator; no impact is
declared. Correct the braking direction and pad before docking.}

\clearpage
\section*{Apollo 11 checkpoint scripts / 11-09--10}
\NarrationCard{11-09}{Bring Columbia home}
{All three crew are back aboard Columbia.
Houston has supplied the Earth-return components and entry clock.
Read the signs, compute the magnitude, and measure coast from cutoff.}
{The return pad is accepted.
The engine fires and shuts down on the checked clock.
The simulator advances Columbia onto its nominal Earth-return branch.}
{Return targeting calls for a trim update.
Review the vector magnitude, acceleration, and cutoff-to-entry
interval before we accept the corrected departure.}
{The entry predictor rejects the return pad.
Rehearse a return-burn replan and correct the signed command.
Do not announce a successful Earth return until the later gates pass.}

\NarrationCard{11-10}{Entry corridor and recovery}
{Earth fills the window. Houston's training predictor supplies
horizontal and downward speeds. Check both the angle corridor
and energy gate before the game commits entry.}
{Both entry checks pass.
In the accepted branch, the capsule descends, the parachutes open,
and Columbia reaches the water on the planned game schedule.
Read the successful Apollo 11 ending once all corrected gates pass.}
{The predictor requests a corrected entry pad.
Hold the entry narration while you revise the angle, energy,
or clock result. Recovery is narrated only after acceptance.}
{This pad is not accepted for entry.
A shallow-angle error warns of skip-out; a steep-angle error warns
of excessive deceleration; an energy error raises heating uncertainty.
Use only the warning that fits this sheet, then correct if reserve permits.}

\clearpage
\section*{Apollo 13 checkpoint scripts / 13-01--02}
\NarrationCard{13-01}{The alarm: power the lifeboat}
{Warning lights are appearing after the service-module accident.
Swigert, check the emergency battery endurance.
Haise, prepare Aquarius. Lovell, authorize the transfer only after
the activation and power-down plan has been checked.}
{The transfer plan is accepted.
Odyssey's nonessential loads are shut down, Aquarius comes online,
and the crew moves into the lifeboat on the checked schedule.}
{The transition checklist threatens the battery margin.
Revise the activation and remaining-energy figures before the
transfer plan is accepted. The accident itself is not your penalty.}
{The power-down budget is rejected.
Rehearse the contingency transfer and correct the pad.
Keeping Odyssey at the emergency load indefinitely is not a rescue plan.}

\NarrationCard{13-02}{Restore a free return}
{Houston needs a return maneuver with the spacecraft still docked.
Haise, work the two-axis pad; Lovell, check direction.
Swigert, confirm that the modeled fuel reserve survives the burn.}
{The free-return pad passes the model's checks.
Aquarius performs the accepted correction.
Houston can now plan the remaining rescue around this return branch.}
{Return tracking predicts a correction.
Review the vector magnitude, burn duration, and fuel estimate
before accepting the nominal maneuver.}
{The free-return predictor rejects this command.
Rehearse the targeting replan, preserving the supplied directions.
There is no lunar-landing alternative in this rescue mission.}

\clearpage
\section*{Apollo 13 checkpoint scripts / 13-03--04}
\NarrationCard{13-03}{Three people, one oxygen budget}
{Aquarius must support three people for the remaining trip.
These are full-horizon planning figures from the stated epoch.
Compare both profiles and choose one that protects the oxygen reserve.}
{The conservation profile and budget are accepted.
The ledger shows oxygen protected for the return.
Keep this corrected plan; do not subtract the full horizon a second time.}
{Houston asks for a tighter ration-sheet review.
Correct the time horizon and total crew use before proceeding.
The warning does not create a new oxygen loss.}
{The selected profile cannot protect the modeled oxygen margin.
Replan conservation immediately in the simulator.
A profile that depletes the pool is not accepted just because
its arithmetic is written neatly.}

\NarrationCard{13-04}{A dark, cold electrical plan}
{The cabin can be quiet and dim, but the accepted plan must keep
energy protected. Separate idle power from the operations allowance
and check the whole return horizon.}
{The low-power plan is accepted.
Unneeded loads stay off while the required systems remain in
the game plan. The crew settles into the dim return coast.}
{The switching sheet needs load shedding.
Correct the power and energy budget; keep the reserve protected.
The colder, darker cabin is narration, not a second numerical penalty.}
{The chosen profile predicts depletion before recovery.
Rewrite the electrical plan before continuing.
Do not borrow the separate command-module entry battery to hide the error.}

\clearpage
\section*{Apollo 13 checkpoint scripts / 13-05--06}
\NarrationCard{13-05}{PC+2: shorten the trip}
{The lunar swingby is behind you.
Houston needs PC+2 checked with the supplied mass model.
Use the logarithmic scale carefully: propellant use must reduce mass.}
{The burn budget is accepted.
Aquarius fires through the planned PC+2 interval and shuts down
at the historical schedule anchor. The game's return branch
now continues toward Earth.}
{The duration differs from the accepted model.
Houston requests a throttle-pad correction.
Recheck the logarithm, subtraction, and division before the burn is accepted.}
{The burn budget or exponent fails the targeting check.
Rehearse emergency targeting and rebuild the mass calculation.
Do not command a model in which the spacecraft gains mass by burning fuel.}

\NarrationCard{13-06}{Cooling water versus drinking water}
{Haise has cooling requirements on one sheet and drinking allowance
on another. Combine the rates for the entire return horizon
without spending the water protected for entry and recovery.}
{The reduced-cooling profile passes the game water gate.
Record the accepted total and remaining water.
The crew continues with the card's drinking allowance preserved.}
{The ration sheet requires an uncomfortable reduced-cooling replan.
Correct the rates and remaining water before proceeding.
The narrative does not subtract additional water.}
{The selected profile consumes protected water.
Rehearse emergency ration planning, then correct the cooling choice.
The three-person drinking allowance cannot be erased to make the sheet pass.}

\clearpage
\section*{Apollo 13 checkpoint scripts / 13-07--08}
\NarrationCard{13-07}{Make the square cartridge fit}
{Houston has an adapter concept for the available cartridges.
Before we narrate assembly, count the required spares and determine
when the existing capacity runs out. Round capacity needs upward.}
{The cartridge allocation and assembly timing pass.
Read the adapter-assembly transition scene.
The accepted capacity plan covers this card's stated interval.}
{The inventory sheet needs review.
Recheck the production rate, allocation, or deadline before assembly
is accepted. A wrong minimum integer count is RED, not merely AMBER.}
{The plan fails the allocation or timing gate.
A scrubber gap would allow carbon-dioxide accumulation in the model.
Rehearse the rushed adapter contingency, then correct the plan
before narrating successful assembly.}

\NarrationCard{13-08}{Navigate without the usual computer}
{Lovell has Earth in view as the exercise's attitude reference.
Haise reads the correction pad.
The angle is measured from negative x toward positive y;
check the quadrant before committing.}
{The signed manual-navigation pad is accepted.
Lovell holds the reference while Haise calls the burn interval.
The simulator accepts the correction toward the entry target.}
{The aim predictor flags drift.
Reset the modeled attitude and correct the angle or burn duration.
Do not change the reference axis to make the result look smaller.}
{The predictor rejects the direction or angle.
A wrong quadrant could miss the entry corridor.
Rehearse manual navigation from the reference again and correct
the signed pad before continuing.}

\clearpage
\section*{Apollo 13 checkpoint scripts / 13-09--10}
\NarrationCard{13-09}{Thread the entry corridor}
{Houston's local predictor says the current approach needs adjustment.
Find the nominal correction and the two corridor bounds.
The supplied angle response is negative for a positive correction.}
{The nominal pad targets the accepted entry angle.
The simulator accepts the final correction.
The crew can now prepare Odyssey's reactivation sequence.}
{The duration produces a corridor warning.
Correct the nominal pad and both bounds before proceeding.
Being somewhere inside the toy corridor does not replace the requested target.}
{The correction pad is rejected.
Too little modeled correction warns of a shallow approach;
too much warns of excessive entry loading.
Select the applicable warning and rework the signed duration solution.}

\NarrationCard{13-10}{Wake Odyssey and come home}
{Swigert is ready to power Odyssey back up.
Haise and Lovell check the protected battery and completion clock.
Aquarius must remain attached until the accepted reactivation is complete.}
{The energy and timing gates pass.
Odyssey is ready before Aquarius is released.
In the accepted branch, the capsule enters, parachutes deploy,
and the crew reaches the water. Read the successful rescue ending.}
{The transfer sequence needs correction.
Hold the jettison narration while you revise the energy or timing sheet.
No successful recovery is announced until the gates are accepted.}
{The recovery predictor rejects the sequence or battery plan.
Premature jettison or insufficient entry energy would compromise
the modeled return. Correct the plan if reserve remains;
otherwise use the failed-rescue simulation ending.}

\clearpage
\section*{Mission endings: read only the earned branch}
\textbf{Apollo 11 / all ten corrected gates; \(R\geq0\).}
``Columbia is on the water. Recovery crews reach the capsule and
the flight crew is brought aboard. Eagle landed, and the crew came
home. Your mission result is crew recovered; landing achieved.''
Then read exactly one science result:
\begin{itemize}
\item No 11-07 warning: ``The record shows full planned science.''
\item 11-07 AMBER: ``An early turnaround reduced the science return;
the record keeps that warning despite the corrected budget.''
\item 11-07 RED: ``The EVA was curtailed to protect the return margin;
the record does not claim the full science schedule.''
\end{itemize}

\textbf{Apollo 11 / \(R<0\) on 11-01 or 11-02.}
``Command reserve is exhausted at the departure or avoidance checkpoint.
Stop the scored run. The training scenario takes its safe Earth-orbit
abort ending; the lunar landing is canceled.
This is a narrated game branch, not a calculated return trajectory.''

\textbf{Apollo 11 / \(R<0\) on 11-03 through 11-06.}
``Command reserve is exhausted before an accepted landing.
Houston cancels the descent mission. The training scenario takes
its safe-return contingency ending; Eagle does not land.
No later calculation can convert this scored abort into a completed mission.''
For 11-05 or 11-06, optionally add:
``The scenario narrates Eagle returning to Columbia rather than
continuing the landing attempt.''

\textbf{Apollo 11 / \(R<0\) on 11-07 through 11-10.}
``The simulation stops here with command reserve exhausted.
We have not established a completed rendezvous or return, as applicable
to this checkpoint. Record simulation terminated; return or rendezvous
unresolved. Do not narrate an automatic rescue of a landed crew.''

\textbf{Apollo 13 / all ten corrected gates; \(R\geq0\).}
``Odyssey is on the Pacific. Recovery personnel reach the capsule,
and Lovell, Haise, and Swigert are brought to safety.
There was no lunar landing, but the rescue mission achieved its goal:
all three crew recovered. Record the remaining command reserve
as your team's command-quality result, not a survival probability.''

\textbf{Apollo 13 / \(R<0\) at any checkpoint.}
``Command reserve is exhausted. Stop the rescue simulation at this
checkpoint; the game has not established crew survival.
Record rescue simulation failed; crew survival not established.
This result does not change history: the actual Apollo 13 crew returned safely.''
